\section{Localization to actual alternating itineraries}
\label{sec:geometry}
The smallest distance between distinct disks is $g_R$, attained only
by the six normal chords between nearest neighbors. The next center
distance is $\sqrt3$. An outgoing ray cannot return to the same convex
disk without another reflection.

\begin{lemma}[Uniform localization]\label{lem:localization}
There are a transverse coordinate radius $a_0>0$ and an excess
$\varepsilon_0>0$, independent of $j\ge1$ and $R\in K$, with the
following property. Every regular $j$-flight segment satisfying
$S_j\tau_R<jg_R+\varepsilon_0$ alternates between one nearest-neighbor
pair. Its impacts lie in the facing arcs $|y_i|<a_0$ about the normal
chord. Every segment in these arcs which satisfies the local specular
stationarity equations is an actual, unobstructed billiard segment.
\end{lemma}
\begin{proof}
Every flight has length at least $g_R$. If their sum has excess less
than $\varepsilon_0$, each individual length is at most
$g_R+\varepsilon_0$. Its target center $e$ satisfies
$|e|\le2R+g_R+\varepsilon_0=1+\varepsilon_0$; taking
$\varepsilon_0<\sqrt3-1$ forces $|e|=1$.

For nearest disks, the minimum of the distance between their boundary
points is unique. By compactness in $R\in K$, chords whose length is
within $\varepsilon_0$ of $g_R$ have both endpoints and direction
uniformly close to that normal chord. This compactness statement can
also be read from the endpoint cosine bound
\[
 n_0\cdot v,\ -n_1\cdot v
 \ge\frac{1-2R-\tau^2}{2R\tau},
\]
whose right side tends to one as $\tau\downarrow g_R$, uniformly on
$K$. Reflection at the target therefore sends the direction close to
$-e$. Distinct nearest-neighbor directions are separated by $\pi/3$.
Taking the neighborhoods sufficiently small forces the next short
flight to go back to the preceding disk. This argument uses the same
neighborhood at every step and hence does not deteriorate with $j$.

The closed normal chord between the pair has positive distance from
all other disks, uniformly on $K$. For the pair $(0,(1,0))$, the other
nearest centers have vertical coordinate $\sqrt3/2$, while
$R\le R_+<1/2$; more distant centers are farther away. Small facing
arcs preserve this separation. Their connecting chords are outgoing
and incoming with normal cosines bounded away from zero. Convexity
prevents a second crossing of the originating disk before the target.
Thus these are first physical hits, and the stationarity equations,
which express the equality of tangential velocity components, are
exactly the specular reflection law. There is no grazing cutoff on a
preselected preparation: the localization follows from the complete
event's excess constraint.
\end{proof}

Fix the pair $(0,(1,0))$ and use a global transverse coordinate $y$.
The facing boundary points are
$(\sqrt{R^2-y^2},y)$ on the first disk and
$(1-\sqrt{R^2-y^2},y)$ on the second. The same symmetric length function
applies in both directions:
\begin{equation}\label{eq:ell}
 \ell_R(u,v)=
 \sqrt{\bigl(1-\sqrt{R^2-u^2}-\sqrt{R^2-v^2}\bigr)^2+(v-u)^2}.
\end{equation}
For endpoints $u=y_0$, $v=y_j$, define the local broken-length action
\begin{equation}\label{eq:action}
 \mathcal L_{j,R}(y_0,\ldots,y_j)=\sum_{i=0}^{j-1}\ell_R(y_i,y_{i+1}).
\end{equation}
The interior variables satisfy $\partial_{y_i}\mathcal L_{j,R}=0$.
We write $W_{j,R}(u,v)$ for the stationary value, whose existence,
uniqueness and uniform estimates are proved next.

The length has the even analytic expansion
\begin{equation}\label{eq:one-quadratic}
 \ell_R(u,v)=g_R+
 \frac{c_R(u^2+v^2)-2uv}{2g_R}+O((u^2+v^2)^2).
\end{equation}
All derivatives of the remainder have uniform analytic bounds in a
fixed small neighborhood for $R\in K$. In particular,
$-\partial_{uv}\ell_R(0,0)=1/g_R$. There is also the elementary bound
\begin{equation}\label{eq:energy-coercive}
 \ell_R(u,v)-g_R\ge
 R-\sqrt{R^2-u^2}+R-\sqrt{R^2-v^2}
 \ge\frac{u^2+v^2}{2R}.
\end{equation}
Consequently any localized orbit of total excess $\varepsilon$ satisfies
$y_0^2+y_j^2+2\sum_{i=1}^{j-1}y_i^2\le2R\varepsilon$.
This bounds all interior variables at once, without a factor $j$.
